\paragraph{Bernoulli\texorpdfstring{$(p)$}{(p)} 基线的归并闭式链：Parry 归一化、代数压力与端点指数}
以下结果在段落 \ref{par:fold-gauge-anomaly-bernoulli-p} 的同一换能器口径下整理，符号与
定理 \ref{thm:fold-gauge-anomaly-bernoulli-p-laws}、
命题 \ref{prop:fold-gauge-anomaly-bernoulli-p-bitpair-law}、
命题 \ref{prop:fold-gauge-anomaly-bernoulli-p-spectrum}、
命题 \ref{prop:fold-gauge-anomaly-bernoulli-p-endpoint-ldp}
保持一致。

\begin{theorem}[规范差压力与 Parry 归一化链的显式闭式]\label{thm:fold-bernoulli-p-parry-pressure-chain}
取 \(p\in(0,1)\)，记 \(q:=1-p\)。令失配计数
\[
G_m=\sum_{t=0}^{m-1}\mathbf 1_{\{X_t\neq Y_t\}},
\]
并定义扭曲矩阵
\begin{equation}\label{eq:fold-bernoulli-p-parry-Atheta}
A_{\theta}(p)=
\begin{pmatrix}
q&q e^\theta&0&p\\
0&0&p e^\theta&0\\
p e^\theta&1&0&0\\
q&0&0&0
\end{pmatrix}.
\end{equation}
记 \(\rho(\theta;p):=\rho(A_\theta(p))\)。则极限压力存在且
\begin{equation}\label{eq:fold-bernoulli-p-parry-pressure}
P_G(\theta;p):=\lim_{m\to\infty}\frac1m\log\EE_p[e^{\theta G_m}]
=\log \rho(\theta;p).
\end{equation}
在 \(\theta=0\) 处有 \(\rho(0;p)=1\)，并可取 Perron 右、左特征向量为
\begin{equation}\label{eq:fold-bernoulli-p-parry-eigenvectors}
r(p)=\bigl(q,\ p^2,\ p,\ q^2\bigr)^{\mathsf T},
\qquad
\ell(p)=\bigl(1,\ 1,\ p,\ p\bigr).
\end{equation}
由 Parry 归一化得到
\begin{equation}\label{eq:fold-bernoulli-p-parry-kernel}
P(p)=\mathrm{diag}(r(p))^{-1}A_0(p)\,\mathrm{diag}(r(p))
=
\begin{pmatrix}
q&p^2&0&pq\\
0&0&1&0\\
q&p&0&0\\
1&0&0&0
\end{pmatrix},
\end{equation}
其平稳分布
\begin{equation}\label{eq:fold-bernoulli-p-parry-stationary}
\pi(p)=\frac{\ell\odot r}{\langle \ell,r\rangle}
=\frac1{1+p^3}\bigl(q,\ p^2,\ p^2,\ p q^2\bigr).
\end{equation}
\end{theorem}
\begin{proof}
\eqref{eq:fold-bernoulli-p-parry-Atheta} 与
\eqref{eq:fold-bernoulli-p-parry-pressure} 即
定理 \ref{thm:fold-gauge-anomaly-bernoulli-p-laws} 的矩阵接口重写。
直接代入可验
\[
A_0(p)\,r(p)=r(p),\qquad
\ell(p)A_0(p)=\ell(p),
\]
得 \eqref{eq:fold-bernoulli-p-parry-eigenvectors}。
由公式
\(
P_{ij}=A_{0,ij}r_j/r_i
\)
得 \eqref{eq:fold-bernoulli-p-parry-kernel}。
再由
\(
\pi_i\propto \ell_i r_i
\)
并归一化，得到 \eqref{eq:fold-bernoulli-p-parry-stationary}。证毕。
\end{proof}

\begin{theorem}[平稳联合律、失配密度与输出压缩闭式]\label{thm:fold-bernoulli-p-bitpair-gamma-q-closed}
在定理 \ref{thm:fold-bernoulli-p-parry-pressure-chain} 的平稳链上，\((X,Y)\) 的单点联合分布为
\begin{align}
\PP_p(X=0,Y=0)&=\frac{1-p-p^2+2p^3-p^4}{1+p^3},\nonumber\\
\PP_p(X=0,Y=1)&=\frac{p^2(1-p)}{1+p^3},\nonumber\\
\PP_p(X=1,Y=0)&=\frac{p^2(2-p)}{1+p^3},\label{eq:fold-bernoulli-p-bitpair-joint-law}\\
\PP_p(X=1,Y=1)&=\frac{p\bigl(p^3+(1-p)^2\bigr)}{1+p^3}. \nonumber
\end{align}
从而失配密度
\begin{equation}\label{eq:fold-bernoulli-p-gamma}
\gamma(p):=\lim_{m\to\infty}\frac1m\EE_p[G_m]
=\PP_p(X\neq Y)
=\frac{p^2(3-2p)}{1+p^3}.
\end{equation}
输出端“1”密度满足
\begin{equation}\label{eq:fold-bernoulli-p-output-density}
q(p):=\PP_p(Y=1)
=\frac{p(p^3-p+1)}{1+p^3},
\qquad
0<q(p)<p\quad(p\in(0,1)).
\end{equation}
\end{theorem}
\begin{proof}
\eqref{eq:fold-bernoulli-p-bitpair-joint-law} 与
\eqref{eq:fold-bernoulli-p-gamma}
由命题 \ref{prop:fold-gauge-anomaly-bernoulli-p-bitpair-law} 直接得到。
\eqref{eq:fold-bernoulli-p-output-density} 的第一式来自
\(
\PP(Y=1)=\PP(0,1)+\PP(1,1)
\)；
第二式由恒等式
\[
p-q(p)=\frac{p^2}{1+p^3}>0
\]
给出。证毕。
\end{proof}

\begin{theorem}[代数压力四次方程与 CLT 方差率闭式]\label{thm:fold-bernoulli-p-pressure-quartic-clt-variance}
令 \(u=e^\theta\)，记 \(\lambda(u,p):=\rho(A_{\log u}(p))\)。
则 \(\lambda(u,p)\) 是下列四次方程在 \(\lambda>0\) 上的唯一最大实根：
\begin{equation}\label{eq:fold-bernoulli-p-pressure-quartic-restated}
\lambda^4+(p-1)\lambda^3+\bigl(p^2-p(u+1)\bigr)\lambda^2
+\bigl(p^3u^3-p^2u^3-p^2u+pu\bigr)\lambda
+\bigl(-p^3u+p^2u\bigr)=0.
\end{equation}
并且 \(P_G(\theta;p)=\log\lambda(e^\theta,p)\) 在 \(\theta\in\RR\) 上解析。

进一步有中心极限定理
\begin{equation}\label{eq:fold-bernoulli-p-clt-restated}
\frac{G_m-m\gamma(p)}{\sqrt m}\ \Rightarrow\ \mathcal N\!\bigl(0,\sigma^2(p)\bigr),
\end{equation}
其中
\begin{equation}\label{eq:fold-bernoulli-p-variance-closed}
\sigma^2(p)=
\frac{p^2(1-p)\bigl(21 p^5-6 p^4+14 p^3-36 p^2+7 p+9\bigr)}
{(p+1)^3(p^2-p+1)^3}.
\end{equation}
\end{theorem}
\begin{proof}
\eqref{eq:fold-bernoulli-p-pressure-quartic-restated} 与
\eqref{eq:fold-bernoulli-p-variance-closed}
与定理 \ref{thm:fold-gauge-anomaly-bernoulli-p-laws} 的四次特征方程与方差率闭式完全等价；
\eqref{eq:fold-bernoulli-p-clt-restated} 则由同一定理的 CLT 结论给出。证毕。
\end{proof}

\begin{proposition}[偏置拆解 Jordan 共振与协方差纯指数混合]\label{prop:fold-bernoulli-p-jordan-resonance-only-half}
矩阵 \(P(p)\) 的谱为
\begin{equation}\label{eq:fold-bernoulli-p-spectrum}
\mathrm{Spec}(P(p))
=\left\{1,\ -p,\ \pm\sqrt{p(1-p)}\right\}.
\end{equation}
仅当 \(p=\tfrac12\) 时发生非主特征值碰撞
\(
-p=-\sqrt{p(1-p)}
\)，并可能出现非半单 Jordan 块；
该点确有 Jordan 退化（见命题 \ref{prop:fold-gauge-anomaly-bernoulli-p-spectrum}）。

对任意 \(p\in(0,1)\setminus\{\tfrac12\}\)，特征值互异，\(P(p)\) 在 \(\RR\) 上可对角化。
因此对任意平方可积边观测
\(
H_t=H(S_t,S_{t+1})
\)
（平稳初值）存在常数 \(C_H(p)<\infty\) 使
\begin{equation}\label{eq:fold-bernoulli-p-covariance-exp-bound}
\bigl|\mathrm{Cov}(H_0,H_k)\bigr|
\le C_H(p)\,\Lambda(p)^k,\qquad
\Lambda(p):=\max\!\left\{p,\sqrt{p(1-p)}\right\}<1.
\end{equation}
特别地，对 \(g_t=\mathbf 1_{\{X_t\neq Y_t\}}\)，当 \(p\neq\tfrac12\) 时协方差无任何 \(k\) 多项式前因子。
\end{proposition}
\begin{proof}
\eqref{eq:fold-bernoulli-p-spectrum}
由命题 \ref{prop:fold-gauge-anomaly-bernoulli-p-spectrum} 直接成立。
当 \(p\neq\tfrac12\) 时可对角化，谱分解给出
\eqref{eq:fold-bernoulli-p-covariance-exp-bound}。
当 \(p=\tfrac12\) 时，\(-\tfrac12\) 二重根的几何重数为 \(1\)，故出现 Jordan 修正。证毕。
\end{proof}

\begin{theorem}[规范差密度的大偏差端点指数]\label{thm:fold-bernoulli-p-endpoint-ldp-restated}
定义
\[
I_p(s):=\sup_{\theta\in\RR}\bigl(s\theta-P_G(\theta;p)\bigr).
\]
则端点指数具有闭式
\begin{equation}\label{eq:fold-bernoulli-p-I1}
I_p(1)=\frac13\log\frac{1}{p^2(1-p)},
\end{equation}
\begin{equation}\label{eq:fold-bernoulli-p-I0}
I_p(0)=
-\log\!\left(\frac{(1-p)+\sqrt{(1-p)(1+3p)}}{2}\right).
\end{equation}
等价地，
\begin{equation}\label{eq:fold-bernoulli-p-endpoint-prob}
\PP_p(G_m=m)=\exp\bigl(-mI_p(1)+o(m)\bigr),\qquad
\PP_p(G_m=0)=\exp\bigl(-mI_p(0)+o(m)\bigr).
\end{equation}
\end{theorem}
\begin{proof}
\eqref{eq:fold-bernoulli-p-I1}--\eqref{eq:fold-bernoulli-p-I0}
由命题 \ref{prop:fold-gauge-anomaly-bernoulli-p-endpoint-ldp} 直接给出；
\eqref{eq:fold-bernoulli-p-endpoint-prob} 是端点点质量的指数等价表述。证毕。
\end{proof}

\begin{proposition}[失配密度的全局最大化与平方最大值恒等式]\label{prop:fold-bernoulli-p-gamma-global-max}
函数
\[
\gamma(p)=\frac{p^2(3-2p)}{1+p^3}\qquad(p\in(0,1))
\]
具有唯一全局极大点 \(p_\star\in(0,1)\)，其由
\begin{equation}\label{eq:fold-bernoulli-p-gamma-root}
p_\star^3+2p_\star-2=0
\end{equation}
唯一刻画。
并且极大值满足精确恒等式
\begin{equation}\label{eq:fold-bernoulli-p-gamma-max-square}
\max_{p\in(0,1)}\gamma(p)=\gamma(p_\star)=p_\star^2.
\end{equation}
\end{proposition}
\begin{proof}
命题
\ref{prop:fold-gauge-anomaly-bernoulli-p-cumulant-denominator-law}
之后的推论 \ref{cor:fold-gauge-anomaly-bernoulli-p-density-unimodal}
已给出 \eqref{eq:fold-bernoulli-p-gamma-root} 与
\eqref{eq:fold-bernoulli-p-gamma-max-square}。证毕。
\end{proof}

\begin{theorem}[Fold 输出密度的偏置衰减、饱和值与唯一拐点]\label{thm:fold-bernoulli-p-output-density-concavity-switch}
令
\[
q(p):=\PP_p(Y=1)=\frac{p(p^3-p+1)}{1+p^3},\qquad p\in[0,1].
\]
则成立：
\begin{enumerate}
  \item 严格单调：\(q\) 在 \([0,1]\) 上严格递增；
  \item 偏置衰减与饱和界：
  \[
  0<q(p)<p\quad(p\in(0,1)),\qquad
  q(p)\le \frac12\quad(p\in[0,1]),
  \]
  且 \(q(p)=\tfrac12\) 当且仅当 \(p=1\)；
  \item 唯一拐点：存在唯一 \(p_{\mathrm{inf}}\in(0,1)\) 使
  \begin{equation}\label{eq:fold-bernoulli-p-output-inflection}
  p_{\mathrm{inf}}^3=\frac{7-3\sqrt5}{2},
  \end{equation}
  并且 \(p<p_{\mathrm{inf}}\) 时 \(q\) 严格凹，\(p>p_{\mathrm{inf}}\) 时 \(q\) 严格凸。
\end{enumerate}
\end{theorem}
\begin{proof}
由有理函数微分得
\[
q'(p)=\frac{p^6+p^4+2p^3-2p+1}{(p+1)^2(p^2-p+1)^2}>0,
\]
故严格递增。

再由
\[
p-q(p)=\frac{p^2}{1+p^3}>0\quad(p\in(0,1))
\]
得 \(0<q(p)<p\)。
并且
\[
q(p)\le \frac12
\iff
2p(p^3-p+1)\le 1+p^3
\iff
(p-1)(2p^3+p^2-p+1)\le 0.
\]
由于 \(2p^3+p^2-p+1>0\) 在 \([0,1]\) 上恒成立，得 \(q(p)\le \tfrac12\)，且仅 \(p=1\) 取等。

二阶导数
\[
q''(p)=
-\frac{2\,(p^6-7p^3+1)}{(p+1)^3(p^2-p+1)^3}.
\]
令 \(t=p^3\)，则符号由 \(t^2-7t+1\) 决定。
在 \((0,1)\) 内仅有小根
\(
t_\ast=\frac{7-3\sqrt5}{2}
\)，对应唯一拐点
\eqref{eq:fold-bernoulli-p-output-inflection}，并给出凹凸区间。证毕。
\end{proof}

\begin{remark}[公开索引同形排查说明]\label{rem:fold-bernoulli-p-closed-form-public-index-note}
围绕闭式表达
\(
\frac{p^2(3-2p)}{1+p^3}
\)、
\eqref{eq:fold-bernoulli-p-variance-closed}、
\eqref{eq:fold-bernoulli-p-I1}--\eqref{eq:fold-bernoulli-p-I0}
及
\eqref{eq:fold-bernoulli-p-output-inflection}
进行公开关键词检索时，未见与本段“同一换能器语境 + 同一参数化 + 同阶同系数闭式链”完全同型的可见条目。
该结论仅是公开索引层面的同形排查，不替代完备文献证明。
\end{remark}

\endinput

